 \chapter{傅里叶级数}\label{ch:b2:fourier}

每一个周期信号能否由纯粹的正弦函数与余弦函数重新构造出来？
傅里叶大胆的“是”开创了一个世纪的分析学。本章
证明了本阶段可及的两块基石：
\emph{狄利克雷定理}（对分段 $C^1$ 函数的逐点重构，
借助狄利克雷核）与\emph{帕塞瓦尔恒等式}（信号的能量
等于其各次谐波能量之和），并
推出若干经典数值级数 --- 其中首推巴塞尔级数 $\sum 1/n^2 =
\pi^2/6$。

全文所涉函数均为 $2\pi$ 周期、分段连续、复值的，
$\mathcal{C}$ 表示其中的连续函数。

\section{傅里叶系数}

\begin{definition}\label{def:b2:fourier:coefficients}
$f$ 的\emph{傅里叶系数}\index{傅里叶系数}
由下式给出：
\[
c_n(f) = \frac{1}{2\pi}\int_{-\pi}^{\pi} f(t)\,\eu^{-\iu n t}\,\dd
t \qquad (n \in \Z),
\]
并记实形式系数 $a_n = c_n + c_{-n}$、$b_n = \iu(c_n
- c_{-n})$，于是\emph{傅里叶部分和}\index{傅里叶
部分和}为
\[
S_N(f)(t) = \sum_{n=-N}^{N} c_n(f)\,\eu^{\iu nt}
= \frac{a_0}{2} + \sum_{n=1}^{N}\bigl(a_n\cos nt + b_n \sin
nt\bigr).
\]
在 $\mathcal{C}$ 上定义埃尔米特内积 $\langle f,
g\rangle = \frac{1}{2\pi}\int_{-\pi}^{\pi}\conj f\,g$，
指数函数 $e_n(t) = \eu^{\iu nt}$ 是\emph{标准正交}的
（$\langle e_m, e_n\rangle = \delta_{mn}$，直接计算），且
$c_n(f) = \langle e_n, f\rangle$：傅里叶分析
就是无穷维的埃尔米特几何（见\cref{ch:b2:hermitian}）。
\end{definition}

\begin{proposition}[贝塞尔不等式]\label{prop:b2:fourier:bessel}
$S_N(f)$ 是 $f$ 在
三角多项式空间 $\mathcal{T}_N$ 上的正交投影（次数 $\leq N$），
且\index{贝塞尔不等式}
\[
\sum_{n=-N}^{N} \abs{c_n(f)}^2 \leq \norm f_2^2 =
\frac{1}{2\pi}\int_{-\pi}^{\pi} \abs f^2 :
\]
级数 $\sum \abs{c_n}^2$ 收敛，且 $c_n(f) \to 0$，
当 $\abs n \to \infty$ 时（系数形式的黎曼–勒贝格引理）。
\end{proposition}

\begin{proof}
$f - S_N(f)$ 与每个 $e_k$ 正交，$\abs k \leq N$（$\langle
e_k, f - S_N f\rangle = c_k - c_k = 0$）：$S_Nf$ 是到
$\mathcal{T}_N = \operatorname{Vect}(e_{-N}, \dots,
e_N)$ 的正交投影（本科一年级分册中的投影定理，
埃尔米特情形下原样成立）。由毕达哥拉斯定理，$\norm f_2^2 = \norm{S_Nf}_2^2 +
\norm{f - S_Nf}_2^2 \geq \norm{S_Nf}_2^2 = \sum_{\abs n \leq N}
\abs{c_n}^2$；令 $N \to \infty$。
\end{proof}

\begin{example}[均方意义下的最佳逼近]\label{ex:b2:fourier:bestapprox}
低次三角多项式逼近
锯齿波 $f(t) = t$（在 $\intoo{-\pi}{\pi}$ 上）的均方逼近效果如何？
由\cref{prop:b2:fourier:bessel}可知，最佳的
$N$ 次逼近\emph{恰为} $S_N(f)$，其平方误差为
\[
\norm{f - S_Nf}_2^2 = \norm f_2^2 -
\sum_{\abs n\leq N}\abs{c_n}^2 .
\]
这里 $\norm f_2^2 = \frac{1}{2\pi}\int_{-\pi}^\pi t^2\dd t =
\frac{\pi^2}{3}$，又由 $b_n = \frac{2(-1)^{n+1}}{n}$
（见\cref{ex:b2:fourier:basel}），$\abs{c_n}^2 + \abs{c_{-n}}^2 =
\frac{b_n^2}{2} = \frac{2}{n^2}$。于是
\[
\norm{f - S_Nf}_2^2
= \frac{\pi^2}{3} - \sum_{n=1}^{N}\frac{2}{n^2}
\qquad\text{：数值上依次为 } 1.29,\ 0.79,\ 0.57,\ 0.44
\]
当 $N = 1, 2, 3, 4$ 时 --- 误差递减，但下降很慢：尾项
$\sum_{n>N}\frac2{n^2} \sim \frac2N$ 由系数的慢速
$\frac1n$ 型衰减所决定，而这本身正是
跳跃的特征（逆用\cref{exo:b2:fourier:6}）。本例
的要点在于：帕塞瓦尔恒等式把逼近的优劣化为某一数值级数的尾项，
并且在看到任何图像之前就预言：跳跃使傅里叶级数
收敛得十分勉强。
\end{example}

\begin{method}[高效计算傅里叶系数]\label{met:b2:fourier:coefficients}
在做任何积分之前，先注意：
\begin{enumerate}
  \item \emph{奇偶性：}偶函数 $f$ 的 $b_n = 0$，奇函数 $f$ 的 $a_n
        = 0$，且剩下的积分化为 $\frac2\pi
        \int_0^\pi$ 形式的积分 --- 工作量减半，可靠性加倍。
  \item \emph{三角多项式本就现成：}
        将乘积线性化（$\cos^3$、$\sin^2\cos$、\dots），并
        直接读出系数
        （见\cref{exo:b2:fourier:9}）；标准正交性使得任何
        进一步的积分均属多余。
  \item \emph{指数因子用复指数处理：}对
        因子 $\eu^{at}$ 或阻尼振荡，直接计算
        $c_n$，只需对 $\eu^{(a - \iu
        n)t}$ 积分一次，胜过两次分部积分
        （见\cref{exo:b2:fourier:10}）。
  \item \emph{对已知展开式求导：}若 $f'$ 的系数已知，
        且 $f$ 连续，则 $c_n(f) =
        \frac{c_n(f')}{\iu n}$（$n \neq 0$）可求得除
        $c_0$ 外的全部系数，此项恰为均值 --- 往往是最快的途径，
        且恰在
        \cref{thm:b2:fourier:parseval}~(1) 的假设下才合法。
\end{enumerate}
\end{method}

\section{狄利克雷定理}

\begin{lemma}[狄利克雷核]\label{lem:b2:fourier:kernel}
$S_N(f)(x) = \frac{1}{2\pi}\int_{-\pi}^{\pi} f(x + u)\,D_N(u)\,\dd
u$，其中\index{狄利克雷核}
\[
D_N(u) = \sum_{n=-N}^{N} \eu^{\iu nu}
= \frac{\sin\bigl((N + \frac12)u\bigr)}{\sin\frac u2}
\quad (u \notin 2\pi\Z),
\qquad
\frac{1}{2\pi}\int_{-\pi}^{\pi} D_N = 1 .
\]
\end{lemma}

\begin{proof}
将 $c_n$ 的定义代入 $S_N$，并交换求和与积分的次序
（这样做是合理的，因为求和为有限和）：
\[
S_N(f)(x)
= \sum_{n=-N}^{N}\Bigl(\frac{1}{2\pi}\int_{-\pi}^{\pi}
f(t)\,\eu^{-\iu nt}\dd t\Bigr)\eu^{\iu nx}
= \frac{1}{2\pi}\int_{-\pi}^{\pi} f(t)
\sum_{n=-N}^{N}\eu^{\iu n(x-t)}\,\dd t ;
\]
作代换 $u = t - x$，并将积分区间平移回
$\intcc{-\pi}{\pi}$，依被积函数的 $2\pi$ 周期性；
对称的指标范围使得 $\sum_n\eu^{-\iu nu} = D_N(u)$。
闭形式由公比为 $\eu^{\iu u}$ 的等比求和给出，
\[
D_N(u) = \eu^{-\iu Nu}\,\frac{\eu^{\iu(2N+1)u} - 1}{\eu^{\iu u} -
1}
= \frac{\eu^{\iu(N + \frac12)u} - \eu^{-\iu(N+\frac12)u}}
{\eu^{\iu u/2} - \eu^{-\iu u/2}} ,
\]
即为正弦之商。其均值为 $1$，仅
$n = 0$ 一项有贡献。
\end{proof}

\begin{theorem}[黎曼–勒贝格引理]\label{thm:b2:fourier:riemannlebesgue}
设 $g$ 在某闭区间上分段连续，
则 $\int_a^b g(t)\sin(\lambda t + \varphi)\,\dd t \to 0$，当 $\lambda
\to +\infty$ 时。\index{黎曼–勒贝格引理}
\end{theorem}

\begin{proof}
用阶梯函数一致逼近 $g$
（不必用到 \cref{thm:b2:funcseq:weierstrass}，
分段连续函数的初等阶梯函数逼近即可
），并对每个阶梯函数显式积分：每一段
贡献 $O\bigl(\frac1\lambda\bigr)$，而逼近误差
贡献 $\varepsilon(b - a)$。此论证已
作为本科一年级分册积分一章的末尾练习完整完成；
对 $C^1$ 分段，可改用分部积分，
并以 $\frac C\lambda$ 为界。
\end{proof}

\begin{example}[系数衰减得有多快？]\label{ex:b2:fourier:decaytable}
黎曼–勒贝格引理说系数趋于 $0$；其
\emph{衰减速率}即光滑性的度量。取自本章及其练习的三个
例子：
\[
\text{方波： } b_n = \frac{4}{\pi n}\ (n\ \text{为奇数}),
\qquad
\abs t : a_n = \frac{-4}{\pi n^2}\ (n\ \text{为奇数}),
\qquad
\abs{\sin t} : a_{2k} = \frac{-4}{\pi(4k^2-1)} .
\]
$f$ 的跳跃（方波、锯齿波）使系数具有
$\frac1n$ 阶：无正规收敛，跳点处出现
吉布斯现象。连续但带尖角（仅 $f'$ 有跳跃）——
则阶提高到 $\frac{1}{n^2}$：正规收敛，
一致重构。一般地，$k$ 阶导数给出 $c_n =
O(n^{-k})$（见\cref{exo:b2:fourier:6}），反之，若谱
衰减快于任何幂次，则 $f$ 必为 $C^\infty$
（此时可合法地逐项求导）。最后的要点：
信号的正则性与谱的衰减是同一信息——
工程师从另一者的斜率读出一者，
而无需画出函数本身。
\end{example}

\begin{theorem}[狄利克雷]\label{thm:b2:fourier:dirichlet}
设 $f$ 是 $2\pi$ 周期且分段 $C^1$ 的函数，则对每个 $x$，
\[
S_N(f)(x) \xrightarrow[N \to \infty]{}
\frac{f(x^+) + f(x^-)}{2}
\]
（即单侧极限的平均值）——特别地，在每个连续点处 $S_N(f)(x) \to
f(x)$。\index{狄利克雷定理}
\end{theorem}


\begin{proof}
由核引理及其均值为 1 的性质，将积分拆为
$u > 0$ 与 $u < 0$ 两部分（各自均值为 $\frac12$）：
\begin{align*}
S_N(f)(x) - \frac{f(x^+) + f(x^-)}{2}
&= \frac{1}{2\pi}\int_{0}^{\pi} \bigl(f(x+u) -
f(x^+)\bigr)D_N(u)\,\dd u\\
&\quad+ \frac{1}{2\pi}\int_{-\pi}^{0}\bigl(f(x+u) -
f(x^-)\bigr)D_N(u)\,\dd u .
\end{align*}
先处理第一项（第二项对称）。记为
\[
\bigl(f(x + u) - f(x^+)\bigr)\,D_N(u)
= \underbrace{\frac{f(x+u) - f(x^+)}{\sin\frac u2}}_{g(u)}\,
\sin\Bigl(\Bigl(N + \frac12\Bigr)u\Bigr) .
\]
函数 $g$ 在 $\intoc{0}{\pi}$ 上分段连续，
并且在 $0^+$ 处有\emph{有限极限}，写作
\[
g(u) = \frac{f(x+u) - f(x^+)}{u}\cdot\frac{u}{\sin\frac u2} ,
\]
第一个因子趋于 $f'(x^+)$（单侧
可导性，由分段 $C^1$ 得出），第二个因子趋于 $2$
（标准极限 $\frac{\sin v}{v} \to 1$ 在 $v = \frac
u2$）：$g(0^+) = 2f'(x^+)$ 存在。于是 $g$ 可分段
连续地延拓到 $\intcc{0}{\pi}$ 上，从而由黎曼–勒贝格引理
（见\cref{thm:b2:fourier:riemannlebesgue}）使积分趋于
$0$。这正是该假设的意义所在：若没有
单侧导数，因子 $\frac{1}{\sin(u/2)}$ 在
$0$ 处趋于无穷，其速度快于黎曼–勒贝格引理所能抵消的程度，而
逐点收敛对仅为连续的
$f$ 确实可能失效——这一缺口由费耶尔定理（周末问题）
借助求平均予以弥补。
\end{proof}

\begin{example}[狄利克雷在跳跃点处的值]\label{ex:b2:fourier:jumpvalue}
取锯齿函数 $f(t) = t$，定义在 $\intoo{-\pi}{\pi}$ 上
（见\cref{ex:b2:fourier:basel}），其周期延拓
在 $t = \pi$ 处由 $f(\pi^-) = \pi$ 跳至 $f(\pi^+) = -\pi$。
狄利克雷定理保证该点的值为 $\frac{\pi + (-\pi)}{2} = 0$，
事实上 $\sum
\frac{2(-1)^{n+1}}{n}\sin nt$ 的每一项在 $t = \pi$ 处都为零，
级数收敛到中点，忽略
两侧的单侧值。若将求值点改为 $t =
\frac\pi2$（连续点）处，则同一级数
变为莱布尼茨的 $\frac\pi4$。同一级数，两种行为——
恰好对应定理的两条结论。
\end{example}

\begin{theorem}[$C^1$ 情形的正规收敛；帕塞瓦尔]\label{thm:b2:fourier:parseval}
\begin{enumerate}
  \item 若 $f$ 连续、$2\pi$-周期且分段 $C^1$，
        则 $c_n(f') = \iu n\,c_n(f)$，$f$ 的傅里叶级数
        在 $\R$ 上\emph{正规}收敛，且其和为 $f$。
  \item （帕塞瓦尔）对任一分段连续的 $2\pi$-周期
        函数 $f$：\index{帕塞瓦尔恒等式}
        \[
        \frac{1}{2\pi}\int_{-\pi}^{\pi}\abs{f}^2
        = \sum_{n=-\infty}^{\infty} \abs{c_n(f)}^2
        = \frac{\abs{a_0}^2}{4} + \frac12\sum_{n\geq1}
        \bigl(\abs{a_n}^2 + \abs{b_n}^2\bigr).
        \]
        \emph{（此处证明 $f$ 连续且分段 $C^1$ 的情形；
        一般情形在本阶段不加证明地承认。）}
\end{enumerate}
\end{theorem}

\begin{proof}
（1）在每个 $C^1$ 段上分部积分（边界项
由连续性与周期性相消）：$c_n(f') = \iu n c_n(f)$。再由
柯西–施瓦茨不等式作用于两个平方可和族
（对 $f'$ 用\cref{prop:b2:fourier:bessel}）：
\[
\sum_{n \neq 0} \abs{c_n(f)} = \sum_{n\neq0}
\frac{\abs{c_n(f')}}{\abs n}
\leq \Bigl(\sum \abs{c_n(f')}^2\Bigr)^{1/2}
\Bigl(\sum_{n\neq0}\frac{1}{n^2}\Bigr)^{1/2} < \infty :
\]
傅里叶级数正规收敛。其和连续，
并由狄利克雷定理在每点都与 $f$ 相等
（由 $f$ 连续，见\cref{thm:b2:fourier:dirichlet}）：级数
一致收敛到 $f$。

（2）对这样的 $f$：$S_N f \to f$ 一致收敛，故 $\norm{f - S_Nf}_2
\leq \norm{f - S_Nf}_\infty \to 0$，由毕达哥拉斯恒等式
($\norm f_2^2 = \sum_{\abs n \leq N}\abs{c_n}^2 + \norm{f -
S_Nf}_2^2$)在取极限后仍成立。实数形式只是
对 $a_n, b_n$ 的记账。
\end{proof}

\begin{example}[$C^1$ 尾项估计的定量化]\label{ex:b2:fourier:tailbound}
在\cref{thm:b2:fourier:parseval}~(1) 的证明中隐藏着一个可用的
估计。对连续、分段 $C^1$ 的 $f$，同样的
柯西–施瓦茨不等式只用于尾项，得到
\[
\sum_{\abs n > N}\abs{c_n(f)}
= \sum_{\abs n > N}\frac{\abs{c_n(f')}}{\abs n}
\leq \Bigl(\sum_{\abs n > N}\abs{c_n(f')}^2\Bigr)^{\!1/2}
\Bigl(\sum_{\abs n>N}\frac{1}{n^2}\Bigr)^{\!1/2}
\leq \norm{f'}_2\,\sqrt{\frac{2}{N}} ,
\]
用到对 $f'$ 的贝塞尔不等式，以及 $\sum_{n>N}n^{-2} \leq \frac1N$。于是
部分和的一致误差满足
\[
\norm{f - S_Nf}_\infty
\leq \sum_{\abs n>N}\abs{c_n(f)}
\leq \norm{f'}_2\,\sqrt{\frac2N} .
\]
对 $f(t) = \abs t$：$\norm{f'}_2 = 1$（导数为
$\pm1$），因此仅十项就能一致地逼近 $\abs t$，
误差在 $\sqrt{0.2} \approx 0.45$ 之内；取 $N = 10^4$ 时误差在
$0.015$ 之内。小结：一阶导数换来了
$\frac{1}{\sqrt N}$ 的一致收敛速度。
与方波的 $\frac1n$ 系数世界相比（那里根本没有一致收敛），
\cref{ex:b2:fourier:decaytable} 中的对照表
便有了具体的数值。
\end{example}

\begin{example}[巴塞尔问题及相关结果]\label{ex:b2:fourier:basel}
设 $f(t) = t$ 定义在 $\intoo{-\pi}{\pi}$ 上，并做 $2\pi$-周期延拓
（锯齿函数，分段 $C^1$）。计算得 $a_n = 0$（由奇性）以及
\[
b_n = \frac{1}{\pi}\int_{-\pi}^{\pi} t\sin nt\,\dd t
= \frac{2(-1)^{n+1}}{n} .
\]
取 $t = \frac\pi2$，由狄利克雷定理恢复莱布尼茨公式 $\frac\pi4 = 1 -
\frac13 + \frac15 - \dots$；由帕塞瓦尔恒等式得
\[
\frac{1}{2\pi}\int_{-\pi}^{\pi} t^2\,\dd t = \frac{\pi^2}{3}
= \frac12\sum_{n\geq1}\frac{4}{n^2}
\quad\Longrightarrow\quad
\boxed{\;\sum_{n\geq1}\frac{1}{n^2} = \frac{\pi^2}{6}\;}
\]
--- 欧拉的巴塞尔级数，两行即得。函数 $f(t) = t^2$
同样得到 $\sum \frac1{n^4} = \frac{\pi^4}{90}$
（见\cref{exo:b2:fourier:3}）。
\end{example}

\begin{example}[完整展开并带内建检验：$\abs{\sin t}$]\label{ex:b2:fourier:abssin}
函数 $f(t) = \abs{\sin t}$ 连续、偶，
$\pi$-周期（从而是 $2\pi$-周期），分段 $C^1$。
偶性使所有 $b_n$ 为零；$a_0 = \frac1\pi\int_{-\pi}^{\pi}
\abs{\sin t}\dd t = \frac4\pi$；当 $n \geq 1$ 时，
由积化和差公式得
\[
a_n = \frac2\pi\int_0^\pi \sin t\cos nt\,\dd t
= \frac{1}{\pi}\int_0^\pi\bigl(\sin(1+n)t +
\sin(1-n)t\bigr)\dd t
= \frac2\pi\cdot\frac{1 + \cos n\pi}{1 - n^2}
\]
当 $n \neq 1$ 时（$a_1 = 0$ 直接验证）：奇数 $n$ 时为零，且
$a_{2k} = \frac{-4}{\pi(4k^2-1)}$。由
\cref{thm:b2:fourier:parseval}~(1)可知级数正规收敛，
并有
\[
\abs{\sin t} = \frac{2}{\pi} - \frac{4}{\pi}
\sum_{k\geq1}\frac{\cos(2kt)}{4k^2 - 1}
\qquad (t \in \R) .
\]
内建检验：取 $t = 0$，恒等式要求
$\sum_{k\geq1}\frac{1}{4k^2-1} = \frac12$，
裂项相消即可验证：
\[
\sum_{k\geq1}\frac{1}{4k^2-1}
= \frac12\sum_{k\geq1}\Bigl(\frac{1}{2k-1} -
\frac{1}{2k+1}\Bigr) = \frac12 . \checkmark
\]
小结：$\abs{\sin}$ 的谱只落在
\emph{偶数}频率上——对正弦做整流会使其频率成分加倍，
频率成分，这正是全波整流器会以
$100$ 或 $120$ 赫兹嗡鸣的原因，恰为市电频率的两倍。
\end{example}

\begin{example}[帕塞瓦尔等式作为计算工具]\label{ex:b2:fourier:zetafourodd}
帕塞瓦尔等式可整批地把展开式化为数值级数。
将其用于 $f(t) = \abs t$（\cref{exo:b2:fourier:2}：$a_0 =
\pi$，$a_n = \frac{-4}{\pi n^2}$ 当 $n$ 为奇数时，其余为零）：
\[
\frac{1}{2\pi}\int_{-\pi}^{\pi}t^2\,\dd t = \frac{\pi^2}{3}
= \frac{a_0^2}{4} + \frac12\sum_{n \text{ 奇数}} a_n^2
= \frac{\pi^2}{4} +
\frac{8}{\pi^2}\sum_{n\text{ 奇数}}\frac{1}{n^4} ,
\]
由此得
\[
\sum_{n\text{ 奇数}}\frac{1}{n^4}
= \frac{\pi^2}{8}\Bigl(\frac{\pi^2}{3} -
\frac{\pi^2}{4}\Bigr) = \frac{\pi^4}{96} .
\]
与\cref{exo:b2:fourier:3}对照：将
$\sum\frac1{n^4}$ 拆为奇、偶两部分，得 $\zeta$ 型的
记账 $S = S_{\mathrm{odd}} + \frac{S}{16}$，于是 $S =
\frac{16}{15}\cdot\frac{\pi^4}{96} = \frac{\pi^4}{90}$ ---
恰为那里借助另一函数求得的值。两个
展开式，同一个数：这种相容性正是
帕塞瓦尔等距性的作用。小结：每一个新的傅里叶展开，都是
一台生成级数恒等式的机器；周末问题的
第 II 部分说明了这台机器为何
永不自相矛盾。
\end{example}

\begin{omfigure}
\begin{tikzpicture}
\begin{axis}[omaxis, width=9.2cm, height=5.6cm,
    xmin=-3.6, xmax=3.6, ymin=-1.5, ymax=1.5,
    xtick={-3.1416, 0, 3.1416},
    xticklabels={$-\pi$, $0$, $\pi$}, ytick={-1,1}]
  \addplot[omExo, thick] coordinates {(-3.1416,-1) (0,-1)};
  \addplot[omExo, thick] coordinates {(0,1) (3.1416,1)};
  \addplot[omDef, thick, domain=-3.3:3.3, samples=200]
    {4/pi*sin(deg(x))};
  \addplot[omThm, very thick, domain=-3.3:3.3, samples=400]
    {4/pi*(sin(deg(x)) + sin(deg(3*x))/3 + sin(deg(5*x))/5 +
     sin(deg(7*x))/7 + sin(deg(9*x))/9)};
\end{axis}
\end{tikzpicture}

{\small 方波（灰色）与傅里叶部分和 $S_1$
（蓝色）与 $S_9$（红色）：在每个连续点处均收敛，但
在跳跃点附近持续存在 $\sim 9\%$ 的过冲 --- 即
\emph{吉布斯现象}\index{吉布斯现象}。一致
收敛恰因极限不连续而失效。}
\end{omfigure}

\begin{example}[平移与调制]\label{ex:b2:fourier:shiftrules}
两条只占一行的规则，可由一个展开式生成许多展开式。设 $a
\in \R$，作代换 $s = t - a$：
\[
c_n\bigl(f(\cdot - a)\bigr)
= \frac{1}{2\pi}\int_{-\pi}^{\pi}f(t - a)\eu^{-\iu nt}\dd t
= \eu^{-\iu na}\,c_n(f)
\qquad\text{（平移调制频谱）},
\]
又由定义直接可得
\[
c_n\bigl(\eu^{\iu kt}f\bigr) = c_{n-k}(f)
\qquad\text{（调制平移频谱）}.
\]
算例：将锯齿波平移 $\pi$，记 $g(t) = f(t -
\pi)$，其中 $f(t) = t$，其 $b_n$ 系数为
$(-1)^n\cdot\frac{2(-1)^{n+1}}{n} = -\frac2n$，于是
$g \sim -2\sum\frac{\sin nt}{n}$ 是在
$0$ 而非 $\pi$ 处跳跃的锯齿波的展开式，无需重算积分。小结
：时间平移只改变相位，决不改变振幅
（$\abs{c_n}$ 对平移不变），故能量
（帕塞瓦尔）与收敛类型都只取决于信号的
\emph{形状}，而与计时起点无关。
\end{example}

\begin{remark}[常见误区]
\emph{(i) 三种收敛，三种度量：}逐点
（狄利克雷：需分段 $C^1$，在
每个跳跃点 --- 取\emph{中点}值，绝不取单侧值），一致（需
连续的极限；跨越跳跃时不可能，吉布斯现象即
其明显症状），以及均方收敛（帕塞瓦尔：最
稳健，对个别点不敏感）。总须指明
所断言的是哪一种。\emph{(ii) 默认不逐项求导
：}对
\cref{ex:b2:fourier:basel}的锯齿波级数逐项求导，得 $\sum
2(-1)^{n+1}\cos nt$，其通项甚至不趋于 $0$ ---
函数列一章中的换序定理需要
\emph{导出}级数的一致收敛，而跳跃恰好
破坏了这一点。先保证光滑，再求导
（\cref{exo:b2:fourier:6} 即是对照表）。\emph{(iii)
对称部分和：}狄利克雷定理涉及 $S_N =
\sum_{-N}^{N}$；重排或先对一侧求和，可能
使发散变为收敛，亦可反之。\emph{(iv)
归一化漂移：}各书的约定互不相同
（前置因子 $\frac{1}{2\pi}$ 或 $\frac1\pi$，周期 $2\pi$ 或
$1$）；可靠的不变量是正交归一关系
--- 在所用约定下重算 $\langle e_m, e_n\rangle$，
然后方可信任任何公式。
\end{remark}

\begin{remark}[应用场合]
帕塞瓦尔等式是傅里叶级数 $L^2$ 理论的萌芽：
本科三年级分册补全了这幅图景（指数函数构成
$L^2$ 的希尔伯特基，映射 $f \mapsto (c_n)$ 是一个
双射等距映射）。在本分册中，周末问题
证明了\emph{费耶尔定理} --- 傅里叶级数的切萨罗平均
对每个连续周期
$f$ 一致收敛 --- 使在本阶段不加证明地承认的一般帕塞瓦尔等式
升为定理，导出三角形式的魏尔斯特拉斯定理，并
带来两项惊人的收获：外尔等分布定理与
等周不等式。应用数学每天都在用本章：
信号的谱、振动系统的谐波，
以及快速傅里叶变换（其有限形式见
\cref{exo:b2:hermitian:10}）。
\end{remark}

\begin{remark}[本分册中的视角]
本章与另外三章相互呼应。回顾：埃尔米特
一章提供了几何（标准正交族、投影、贝塞尔不等式），函数列
一章提供了分析（一致收敛、转移定理、
近似恒等式——费耶尔核之于傅里叶级数，
恰如伯恩斯坦多项式之于魏尔斯特拉斯）。侧面看，幂级数
一章的边界理论经阿贝尔求和重新出现，
此处由泊松核实现
（见\cref{exo:b2:fourier:12}）——圆盘半径 $r$ 充当求和参数。
向前看，微分方程一章把周期性强迫分解为谐波，
并将每个谐波送入振子的频率响应；
共振发生于输入的某个傅里叶模式
与固有频率相合之时，
这正是为何该章的周末问题
与本章的周末问题
是同一故事的两半。
\end{remark}

\begin{omfigure}
\begin{tikzpicture}
\begin{axis}[omaxis, width=9.2cm, height=5.6cm,
    xmin=-3.6, xmax=3.6, ymin=-2.5, ymax=9.5,
    xtick={-3.1416, 0, 3.1416},
    xticklabels={$-\pi$, $0$, $\pi$}, ytick={8}]
  \addplot[omDef, thick, domain=0.02:3.3, samples=300]
    {sin(deg(8.5*x))/sin(deg(0.5*x))};
  \addplot[omDef, thick, domain=-3.3:-0.02, samples=300]
    {sin(deg(8.5*x))/sin(deg(0.5*x))};
  \addplot[omThm, very thick, domain=0.02:3.3, samples=300]
    {sin(deg(4*x))^2/(8*sin(deg(0.5*x))^2)};
  \addplot[omThm, very thick, domain=-3.3:-0.02, samples=300]
    {sin(deg(4*x))^2/(8*sin(deg(0.5*x))^2)};
  \node[font=\small, omDef] at (axis cs:1.05,5.2) {$D_8$};
  \node[font=\small, omThm] at (axis cs:-0.75,6.2) {$F_8$};
\end{axis}
\end{tikzpicture}

{\small 狄利克雷核 $D_8$（蓝色）振荡，且取负值；
费耶尔核 $F_8$（红色）非负，
集中于 $0$，均值为 $1$，是一个\emph{近似恒等式}。
非负性恰恰是狄利克雷核所缺乏的，
而正是它使周末问题中的费耶尔定理
成为无条件成立的。}
\end{omfigure}

\section{练习}

\begin{exercise}[$\star$]\label{exo:b2:fourier:1}
计算方波（$f = -1$ 在
$\intoo{-\pi}{0}$ 上，$+1$ 在 $\intoo{0}{\pi}$ 上）的傅里叶系数，
写出狄利克雷在 $t = \frac\pi2$ 及跳跃点 $t = 0$ 处的结论，
并由此导出莱布尼茨级数。
\end{exercise}

\begin{exercise}[$\star$]\label{exo:b2:fourier:2}
将 $f(t) = \abs t$（$\abs t \leq \pi$，$2\pi$ 为周期）
展开为傅里叶级数；说明其级数是\emph{正规}收敛的；在 $t =
0$ 处取值，得 $\sum_{k\geq0}\frac{1}{(2k+1)^2} = \frac{\pi^2}{8}$，
并由此重新推出巴塞尔级数之和。
\end{exercise}

\begin{exercise}[$\star$]\label{exo:b2:fourier:3}
将 $f(t) = t^2$（$\abs t \leq \pi$）展开，并由此推出
\[
\sum_{n\geq1}\frac{(-1)^{n+1}}{n^2} = \frac{\pi^2}{12},
\qquad
\sum_{n\geq1}\frac{1}{n^4} = \frac{\pi^4}{90}
\quad\text{（帕塞瓦尔）}.
\]
\end{exercise}

\begin{exercise}[$\star\star$]\label{exo:b2:fourier:4}
设 $f$ 为连续的 $2\pi$ 周期函数，且 $c_n(f) = 0$，
对一切 $n$ 成立。证明 $f = 0$ \emph{（帕塞瓦尔——本章在哪一类函数上证明了它？
说明连续与分段 $C^1$ 这一条件可以去掉，
方法是承认一般帕塞瓦尔等式，
或概述稠密性论证）}。
\end{exercise}

\begin{exercise}[$\star\star$]\label{exo:b2:fourier:5}
对 $\alpha \notin \Z$，将 $f(t) = \cos(\alpha t)$（$\abs t
\leq \pi$）展开，并由此推出
余切的部分分式展开：
\[
\pi\cot(\pi\alpha) = \frac{1}{\alpha} + \sum_{n\geq1}
\frac{2\alpha}{\alpha^2 - n^2} .
\]
\end{exercise}

\begin{exercise}[$\star\star$]\label{exo:b2:fourier:6}
证明：若 $f$ 为 $2\pi$ 周期的 $C^k$ 函数，且 $f^{(k)}$
分段连续，则 $c_n(f) = O\bigl(\abs n^{-k}\bigr)$：
信号的光滑性 $=$ 其谱的衰减。
\end{exercise}

\begin{exercise}[$\star\star\star$]\label{exo:b2:fourier:7}
（维廷格不等式）设 $f$ 为 $C^1$ 类、$2\pi$ 周期函数，且
$\int_{-\pi}^{\pi} f = 0$。证明
\[
\int_{-\pi}^{\pi} \abs{f}^2 \leq \int_{-\pi}^{\pi} \abs{f'}^2 ,
\]
等号成立当且仅当 $f(t) = a\cos t + b\sin t$。\emph{（在两侧使用帕塞瓦尔
等式；比较 $\abs{c_n}^2$ 与 $n^2\abs{c_n}^2$。）}
\end{exercise}

\begin{exercise}[$\star\star\star$]\label{exo:b2:fourier:8}
（吉布斯常数）对
\cref{exo:b2:fourier:1}中的方波，计算部分和在 $x_N =
\frac{\pi}{2N}$ 处的值：记 $u_k = \frac{(2k+1)\pi}{2N}$，并记 $\Delta u
= \frac{\pi}{N}$，证明
\[
S_{2N-1}\Bigl(\frac{\pi}{2N}\Bigr)
= \frac{4}{\pi}\sum_{k=0}^{N-1} \frac{\sin u_k}{2k+1}
= \frac{2}{\pi}\sum_{k=0}^{N-1} \frac{\sin u_k}{u_k}\,\Delta u
\xrightarrow[N\to\infty]{}
\frac{2}{\pi}\int_0^{\pi}\frac{\sin u}{u}\,\dd u \approx 1.179 :
\]
这是 $\frac{2}{\pi}\cdot\frac{\sin u}{u}$ 在
$\intcc{0}{\pi}$ 上取中点的黎曼和。由此可知，越过
跳跃值 $1$ 的过冲并不随 $N \to \infty$ 而消失。
\end{exercise}

\begin{exercise}[$\star$]\label{exo:b2:fourier:9}
将 $\cos^3 t$ 与 $\sin^2 t\,\cos t$ 展开为傅里叶级数
\emph{（化为线性形式；三角多项式本身即其傅里叶
级数，由系数的唯一性可知）}。求 $c_n$、
$a_n$、$b_n$ 各是多少？
\end{exercise}

\begin{exercise}[$\star\star$]\label{exo:b2:fourier:10}
设 $a > 0$，令 $f(t) = \eu^{at}$ 在 $\intoc{-\pi}{\pi}$ 上，
以 $2\pi$ 为周期作延拓，计算
\[
c_n(f) = \frac{(-1)^n\sinh(a\pi)}{\pi(a - \iu n)} ,
\]
在跳跃点 $t = \pi$ 处应用狄利克雷定理，并
推出双曲余切的部分分式展开：
\[
\coth(\pi a) = \frac{1}{\pi a} + \sum_{n\geq1}
\frac{2a}{\pi(a^2 + n^2)} .
\]
\end{exercise}

\begin{exercise}[$\star\star$]\label{exo:b2:fourier:11}
（卷积）设 $f, g$ 连续且以 $2\pi$ 为周期，定义
\[
(f * g)(x) = \frac{1}{2\pi}\int_{-\pi}^{\pi} f(x - t)\,g(t)\,
\dd t .
\]
证明 $f * g = g * f$，$c_n(f*g) = c_n(f)\,c_n(g)$
\emph{（为交换连续被积函数的两个积分的次序，
比较上限处的两个函数：二者均在
左端点处为零，且导数相同，由连续性及
积分号下求导可得）}，并证明 $S_N(f) =
f * D_N$（狄利克雷核）。（周末问题
中的费耶尔平均同样是卷积，$\sigma_N(f) = f *
F_N$。）
\end{exercise}

\begin{exercise}[$\star\star\star$]\label{exo:b2:fourier:12}
（泊松核：傅里叶级数的阿贝尔平均）。对 $0 \leq r <
1$，令 $P_r(t) = \sum_{n\in\Z} r^{\abs n}\eu^{\iu nt}$。
\begin{enumerate}
  \item 对两个几何级数求和，并证明
        \[
        P_r(t) = \frac{1 - r^2}{1 - 2r\cos t + r^2} > 0,
        \qquad
        \frac{1}{2\pi}\int_{-\pi}^{\pi}P_r = 1 .
        \]
  \item 证明：对 $\delta \leq \abs t \leq \pi$，$P_r(t)
        \leq \frac{1 - r^2}{1 - 2r\cos\delta + r^2} \to 0$，
        $r \to 1^-$ 时，此界一致地趋于零。
  \item 由此可知，对每个连续且以 $2\pi$ 为周期的 $f$，
        其\emph{阿贝尔平均} $(f * P_r)(x) = \sum_n
        r^{\abs n}c_n(f)\,\eu^{\iu nx}$ 一致收敛于 $f$，
        当 $r \to 1^-$ 时 --- 连续参数下的
        费耶尔定理的对应，也是傅里叶级数中的
        阿贝尔求和法的体现，源自幂级数
        那一章。
\end{enumerate}
\end{exercise}

\section{问题：费耶尔定理及其收获}

\begin{problem}\label{pb:b2:fourier:1}
狄利克雷定理要求 $f$ 分段 $C^1$；对仅连续的
$f$，部分和 $S_N(f)$ 可能表现很差。费耶尔
的发现是：其\emph{切萨罗平均}从不如此。其动力是
费耶尔核的正性，而收获极为丰富：
一致三角逼近（魏尔斯特拉斯），
傅里叶系数的唯一性、对本章每个函数的帕塞瓦尔等式
（去掉
\cref{thm:b2:fourier:parseval}中的“承认”），外尔等分布
定理，以及 --- 为几何学的百年历程画上句号的 ---
等周不等式。\index{费耶尔定理}
\index{等分布}\index{等周不等式}
全文中 $f$ 以 $2\pi$ 为周期且分段连续，
并记
\[
\sigma_N(f) = \frac{S_0(f) + S_1(f) + \dots +
S_{N-1}(f)}{N} .
\]

\medskip
\noindent\textbf{第一部分 --- 费耶尔核。}
\begin{enumerate}
  \item 证明：$\sigma_N(f)(x) =
        \frac{1}{2\pi}\int_{-\pi}^{\pi}f(x+u)\,F_N(u)\,\dd u$
        其中 $F_N = \frac{D_0 + \dots + D_{N-1}}{N}$，并
        证明 $\frac{1}{2\pi}\int_{-\pi}^{\pi}F_N = 1$。
  \item 对 $u \notin 2\pi\Z$，证明闭形式：
        \[
        F_N(u) = \frac{1}{N}\,
        \frac{\sin^2\bigl(\frac{Nu}{2}\bigr)}
        {\sin^2\bigl(\frac u2\bigr)} \;\geq\; 0
        \]
        \emph{（对 $\sin\bigl((n+\frac12)u\bigr)$ 求和，视作几何级数的
        虚部）}。
  \item 证明集中性估计：对 $0 < \delta \leq
        \abs u \leq \pi$,
        \[
        F_N(u) \leq \frac{1}{N\sin^2\frac\delta2}
        \xrightarrow[N\to\infty]{} 0
        \quad\text{一致收敛} :
        \]
        $(F_N)$ 是正的近似单位元。
  \item （费耶尔定理）证明：若 $f$ 连续且
        以 $2\pi$ 为周期，则 $\sigma_N(f) \to f$
        \emph{一致收敛}于 $\R$ 上，\emph{（将
        $\bigl(f(x+u) - f(x)\bigr)F_N(u)$ 的积分在 $\abs u =
        \delta$ 处分段；利用海涅定理及第 1--3 小题）}。
  \item 对分段连续的 $f$，证明逐点
        形式 $\sigma_N(f)(x) \to \frac{f(x^+) +
        f(x^-)}{2}$ 在每点 $x$ 处成立，并证明一致有界
        $\norm{\sigma_N(f)}_\infty \leq \norm f_\infty$
        \emph{（正性！）}。
\end{enumerate}

\medskip
\noindent\textbf{第 II 部分 --- 魏尔斯特拉斯，唯一性，
帕塞瓦尔。}
\begin{enumerate}[resume]
  \item （三角形式的魏尔斯特拉斯定理）推出：每个连续的
        $2\pi$ 周期函数
        都是三角多项式的一致极限。
  \item （唯一性）推出：若连续函数 $f$ 满足 $c_n(f) = 0$，
        对所有 $n$ 成立，则它恒为零 --- 两个具有相同
        傅里叶系数的连续周期函数必相等
        （见\cref{exo:b2:fourier:4}，此处不再需要
        任何承认的结论）。
  \item （帕塞瓦尔恒等式，连续情形）利用
        $S_N$ 的投影性质（\cref{prop:b2:fourier:bessel}）及
        $\sigma_Nf \in \mathcal T_{N-1} \subseteq \mathcal
        T_N$，证明
        \[
        \norm{f - S_Nf}_2 \leq \norm{f - \sigma_Nf}_2
        \leq \norm{f - \sigma_Nf}_\infty
        \xrightarrow[N\to\infty]{} 0 ,
        \]
        并由此得到对每个
        \emph{连续}的 $2\pi$ 周期函数 $f$ 成立的帕塞瓦尔恒等式。
  \item （帕塞瓦尔恒等式，分段连续情形）给定 $f$
        分段连续，且 $\varepsilon > 0$，构造一个
        连续周期函数 $g$，使 $\norm{f - g}_2 \leq
        \varepsilon$，\emph{（以仿射插值代替 $f$
        在跳跃点附近的微小区间上的部分）}，
        由此推出 $\norm{f - S_Nf}_2 \to 0$
        \emph{（利用贝塞尔不等式：$\norm{S_Nh}_2 \leq \norm h_2$）}：
        帕塞瓦尔恒等式在
        \cref{thm:b2:fourier:parseval}中所叙述的完全一般情形下成立 --- 此前的“承认”
        已不再需要。
  \item （费耶尔求和无吉布斯现象）与
        \cref{exo:b2:fourier:8} 对比，证明对于方波
        $f$，处处有 $\abs{\sigma_N(f)} \leq 1$，对每个
        $N$ 都成立 --- 切萨罗平均消去了
        困扰 $S_N$ 的过冲。用一句话说明
        $F_N$ 的哪一性质起了作用。
\end{enumerate}

\medskip
\noindent\textbf{第 III 部分 --- 收敛速率。}
\begin{enumerate}[resume]
  \item 证明两个核估计，其中 $0 < \abs u \leq
        \pi$：
        \[
        F_N(u) \leq N,
        \qquad
        F_N(u) \leq \frac{\pi^2}{N u^2}
        \]
        \emph{（对第一式，由归纳法证明 $\abs{\sin N\theta} \leq
        N\abs{\sin\theta}$；对第二式，
        $\sin\frac u2 \geq \frac{u}{\pi}$ 在
        $\intcc{0}{\pi}$ 上成立）}。
  \item 由此推出一阶矩估计
        \[
        \frac{1}{2\pi}\int_{-\pi}^{\pi}\abs u\,F_N(u)\,\dd u
        \;\leq\; \frac{C\,\ln N}{N}
        \qquad (N \geq 2)
        \]
        对某个显式常数 \emph{（在 $\abs u =
        \frac1N$ 处分段）}。
  \item 由此得出：若 $f$ 为 $L$-利普希茨且 $2\pi$ 周期，
        则
        \[
        \norm{\sigma_N f - f}_\infty \leq \frac{C\,L\ln N}{N} .
        \]
  \item （饱和性）计算 $\sigma_N(e_1)$，其中 $e_1(t) =
        \eu^{\iu t}$ 并证明 $\norm{\sigma_Ne_1 - e_1}_\infty
        = \frac1N$：即使对最光滑的函数，费耶尔
        求和的收敛速度也不快于 $\frac1N$ --- 这恰是
        函数列一章周末问题中伯恩斯坦饱和现象的
        精确类比。
  \item （局部性）证明：若分段连续的 $f$
        在 $\intoo{x - \delta}{x + \delta}$ 上恒为零，则
        $\sigma_N(f)(x) \to 0$，无论 $f$ 在别处多么不规则
        --- 在 $x$ 处的均值收敛只依赖 $f$
        在 $x$ 附近的值。
\end{enumerate}

\medskip
\noindent\textbf{第 IV 部分 --- 外尔等分布定理。}
数列 $(x_n)_{n\geq1}$ 位于 $\intco{0}{1}$ 中，称其为
\emph{等分布的}，当对每个区间
$\intcc{a}{b} \subseteq \intcc{0}{1}$，
\[
\frac{\#\{n \leq N : x_n \in \intcc ab\}}{N}
\xrightarrow[N\to\infty]{} b - a .
\]
\begin{enumerate}[resume]
  \item 证明：$(x_n)$ 是等分布的，只要
        $\frac1N\sum_{n\leq N}f(x_n) \to \int_0^1 f$ 对每个
        \emph{连续}的 $1$ 周期函数 $f$ \emph{（将
        $\intcc ab$ 的示性函数夹在两个连续
        分段仿射函数之间，且二者积分相差
        $\varepsilon$）}。
  \item （外尔判别法，充分性）设
        \[
        \frac{1}{N}\sum_{n=1}^{N}\eu^{2\iu\pi kx_n}
        \xrightarrow[N\to\infty]{} 0
        \qquad\text{对每个 } k \in \Z\setminus\{0\} .
        \]
        先证明 $\frac1N\sum f(x_n) \to \int_0^1f$
        对三角多项式成立，再对所有连续的
        $1$ 周期函数 $f$ 成立（由第 6 小题，经平移到周期
        $1$）：结合第 16 小题，$(x_n)$ 是等分布的。
  \item 设 $\alpha$ 为无理数，$x_n = \{n\alpha\}$
        （小数部分）。估计几何和
        \[
        \Bigl|\sum_{n=1}^{N}\eu^{2\iu\pi kn\alpha}\Bigr|
        \leq \frac{2}{\abs{1 - \eu^{2\iu\pi k\alpha}}}
        \qquad (k \neq 0),
        \]
        并由此得出\emph{外尔定理}：$(\{n\alpha\})$ 在
        $\intco{0}{1}$ 中等分布。
  \item 由此推出 $(\{n\alpha\})$ 在
        $\intcc{0}{1}$ 中稠密，其中 $\alpha$ 为无理数，并
        用一句话说明为何等分布性严格强于
        稠密性。
  \item （首位数字）证明：在满足
        $n \leq N$ 的整数中，使 $2^n$ 首位（十进制）数字
        $1$ 的比例趋于 $\log_{10}2 \approx 0.301$ \emph{（首位
        数字为 $1$ 即 $\{n\log_{10}2\} \in
        \intco{0}{\log_{10}2}$；证明 $\log_{10}2$ 为
        无理数）}。
\end{enumerate}

\medskip
\noindent\textbf{第五部分 --- 等周不等式。}设
$\Gamma$ 为闭简单 $C^1$ 曲线，长为 $L$，所围
有向面积为 $A$，以按比例弧长为参数：$z(t) =
x(t) + \iu y(t)$，以 $2\pi$ 为周期，且 $\abs{z'(t)} =
\frac{L}{2\pi}$ 为常数；所围面积为
\[
A = \frac12\int_0^{2\pi}\bigl(x\,y' - y\,x'\bigr)\dd t
= \frac{1}{2}\,\Im\int_0^{2\pi}\conj{z}\,z'\,\dd t
\]
（此处取作有向面积的定义；重积分一章
将证明它与直观面积一致，
借助格林公式）。
\begin{enumerate}[resume]
  \item 将 $z(t) = \sum_{n\in\Z}c_n\eu^{\iu nt}$ 展开（此为
        $C^1$ 函数的正规收敛级数），并
        对 $z'$ 用帕塞瓦尔等式证明：
        \[
        \frac{L^2}{2\pi} = \int_0^{2\pi}\abs{z'}^2\dd t
        = 2\pi\sum_{n\in\Z}n^2\abs{c_n}^2 .
        \]
  \item 同理证明 $A = \pi\sum_{n\in\Z}
        n\,\abs{c_n}^2$ \emph{（极化形式的帕塞瓦尔等式：
        $\frac{1}{2\pi}\int\conj f g = \sum
        \conj{c_n(f)}c_n(g)$，取 $f = z$，$g = z'$）}。
  \item （赫尔维茨）由此得出：
        \[
        L^2 - 4\pi A = 4\pi^2\sum_{n\in\Z}
        (n^2 - n)\abs{c_n}^2 \;\geq\; 0 ,
        \]
        等号成立当且仅当 $z(t) = c_0 +
        c_1\eu^{\iu t}$ --- 即圆。
        等周不等式：在长为 $L$ 的闭曲线中，只有
        圆能围出面积 $\frac{L^2}{4\pi}$。
  \item 验算：对半径为
        $R$ 的圆验证等号成立，对
        边长为 $a$ 的正方形验证严格不等式；说明为何对每个整数都有 $n^2 - n \geq 0$，
        包括负整数在内的 $n$，并指出在何处
        用到了参数化的匀速性。
  \item 综合。每项各用一句话作答：(i) 第一至第三部分赖以推出的
        $F_N$ 所具有的唯一性质，
        而 $D_N$ 所缺少的；(ii) 此处的切萨罗求和
        与幂级数一章的周末问题有何关系；
        （弗罗贝尼乌斯）；(iii) 哪一部分只用到了
        魏尔斯特拉斯（第 6 题），哪一部分需要完整的
        帕塞瓦尔等式；(iv) 用一句话说明本科三年级分册
        增加了什么（$L^2$ 完备性：傅里叶级数作为希尔伯特
        基）。
\end{enumerate}
\end{problem}
